\documentclass[11pt]{article}
\usepackage[margin=1in]{geometry}
\usepackage{amsmath,amssymb}
\usepackage{xcolor}
\usepackage{booktabs}
\usepackage{enumitem}
\usepackage{fancyhdr}
\usepackage{tcolorbox}
\tcbuselibrary{skins,breakable}
\usepackage{hyperref}
\hypersetup{colorlinks=true, linkcolor=blue, urlcolor=blue}

\definecolor{primary}{RGB}{30,64,140}
\definecolor{tipbg}{RGB}{235,244,255}
\definecolor{examplebg}{RGB}{240,248,235}
\definecolor{formulabg}{RGB}{255,247,224}
\definecolor{warnbg}{RGB}{255,236,236}

\pagestyle{fancy}
\fancyhf{}
\renewcommand{\headrulewidth}{0.4pt}
\cfoot{\thepage}

\newtcolorbox{tipbox}[1][Tip]{
  colback=tipbg, colframe=primary, fonttitle=\bfseries,
  title=#1, breakable, rounded corners
}
\newtcolorbox{examplebox}[1][Example]{
  colback=examplebg, colframe=primary!70!black, fonttitle=\bfseries,
  title=#1, breakable, rounded corners
}
\newtcolorbox{formulabox}[1][Formula]{
  colback=formulabg, colframe=primary!70!black, fonttitle=\bfseries,
  title=#1, breakable, rounded corners
}
\newtcolorbox{warnbox}[1][Warning]{
  colback=warnbg, colframe=red!60!black, fonttitle=\bfseries,
  title=#1, breakable, rounded corners
}


\title{\color{primary}\textbf{Time \& Work --- Fundamentals}\\[4pt]\large Work, Rate, Efficiency \& the Core Formulas}
\author{Aptitude Grind}
\date{}

\lhead{\small Time \& Work --- Fundamentals}
\rhead{\small Aptitude Grind}

\begin{document}
\maketitle
\thispagestyle{fancy}

\section*{Why Time \& Work Matters for TCS NQT}
Time \& Work is one of the highest-frequency topics in TCS NQT (Ninja + Digital) Numerical Ability,
alongside Percentages and Ratios. Unlike SSC CGL or CAT, TCS rarely asks difficult algebraic
work-time puzzles --- it sticks to \textbf{direct formula-based questions with at most 1--2 twists}:
ratio based worker efficiency, pipes \& cisterns, alternate working/resting, men-women-children
efficiency, and machine problems. The entire topic rests on one idea: \emph{work, rate, and time are
locked together by a single equation, and almost every question is just that equation applied from a
different angle.}

\section{The Golden Rule}
\begin{formulabox}[Core Identity]
\[
\textbf{Work} = \textbf{Rate} \times \textbf{Time}, \quad \text{i.e. } W = R \times T
\]
\begin{itemize}[leftmargin=1.4em]
  \item \textbf{Work} --- the total task, conventionally treated as \textbf{1 unit} (a whole job)
    unless stated otherwise.
  \item \textbf{Rate} --- work done per unit time (also called \emph{efficiency} or \emph{one
    day's/hour's work}).
  \item \textbf{Time} --- days/hours taken.
\end{itemize}
If total work $= 1$, then $\text{Rate} = 1/\text{Time}$.
\end{formulabox}

This single idea solves almost every question in this topic --- everything below is a variation of
it.

\section{One Day's Work}
If \textbf{A} finishes a work in $N$ days, A's one day's work is $1/N$.

\begin{examplebox}[Example]
A completes a work in 8 days $\Rightarrow$ one day's work $= 1/8$.
\end{examplebox}

The relationship also runs backward: if A's one day's work is $1/15$, then A alone takes
\textbf{15 days} to finish the whole job. Time and one-day's-work are reciprocals of each other.

\section{Combined Work --- Two Workers}
If A finishes a job alone in $x$ days and B finishes it alone in $y$ days, their combined one-day's
work is $1/x + 1/y$, so together they take:
\begin{formulabox}[Combined Work Formula]
\[
\text{Time (A + B)} = \frac{xy}{x+y}
\]
\end{formulabox}
This formula is extremely common and worth memorizing outright.

\begin{examplebox}[Worked Example]
A = 10 days, B = 15 days.
\[
\text{Together} = \frac{10 \times 15}{10 + 15} = \frac{150}{25} = 6 \text{ days}
\]
\end{examplebox}

\section{Combined Work --- Three or More Workers}
For three workers:
\[
\text{Combined one-day's work} = \frac{1}{x} + \frac{1}{y} + \frac{1}{z}
\]
Take the LCM of $x$, $y$, $z$ to add these cleanly (see \texttt{02-fast-tricks-and-shortcuts.tex} for
the LCM shortcut). For $n$ workers, the pattern generalizes --- just sum every individual rate.

\section{The Difference Formula (Finding One Worker From a Pair)}
If A alone takes $x$ days, and A + B together take $y$ days, B's one-day's work is:
\begin{formulabox}[Difference Formula]
\[
\frac{1}{B} = \frac{1}{y} - \frac{1}{x}
\]
\end{formulabox}
then invert to get B's number of days.

\begin{examplebox}[Worked Example]
A = 30 days, A + B = 12 days.
\[
\frac{1}{B} = \frac{1}{12} - \frac{1}{30} = \frac{5}{60} - \frac{2}{60} = \frac{3}{60} = \frac{1}{20}
\Rightarrow B = 20 \text{ days}
\]
\end{examplebox}

This ``given one worker + the pair, find the other worker'' pattern is one of the most frequently
tested variants of Time \& Work.

\section{Efficiency Concept (Very Important)}
\begin{formulabox}[Inverse Proportionality]
\[
\text{Efficiency} \propto \frac{1}{\text{Time}}
\]
Higher efficiency means less time, and the two are always \textbf{inversely proportional}.
\end{formulabox}

If efficiency ratio $A:B = 2:3$, then time ratio $A:B = 3:2$ --- always the reverse of the efficiency
ratio.

\begin{examplebox}[Example]
Time ratio $2:5$ $\Rightarrow$ efficiency ratio $5:2$.
\end{examplebox}

\subsection{Worker-Ratio Questions --- Think in Rates, Not Days}
\textbf{Example:} ``A is twice as efficient as B.''

\begin{tipbox}[Right Approach]
Assume work \emph{rates} directly. Let B's rate $= 1$ unit/day, then A's rate $= 2$ units/day. Never
assume days first --- assume the rate ratio, since that's what the problem is actually telling you.
\end{tipbox}

This single habit --- assigning rates instead of days whenever a ratio is given --- is the fastest
problem-solving upgrade in this entire topic.

\section{Men, Women, and Children Equivalence}
Frequently, a question gives an equivalence like ``2 Men = 3 Women'' and expects you to convert
everyone into one common unit before adding rates.

\begin{examplebox}[Worked Example]
$2M = 3W$ (i.e. equal work done in equal time)
\[
2M = 3W \;\Rightarrow\; M:W = 3:2 \;\Rightarrow\; 1 \text{ man} = \tfrac{3}{2} \text{ women (in work
output)}
\]
\end{examplebox}

Always convert every category (men, women, children) into a single common unit of efficiency before
combining rates --- treat this exactly like the ratio-to-rate trick above.

\section{Machine / Printer / Robot / Server Problems}
Conceptually \textbf{identical} to worker problems --- a machine, printer, robot, or server is just
another ``worker'' with its own rate. No new formula is needed; apply everything above directly.

\section{Pipes and Cisterns}
Same core concept, with a sign convention:
\begin{formulabox}[Pipe Sign Convention]
\centering
\begin{tabular}{lcc}
\toprule
\textbf{Pipe type} & \textbf{Effect} & \textbf{Rate} \\
\midrule
Filling (inlet) & Positive work & $+1/x$ (fills in $x$ hours) \\
Emptying (outlet) & Negative work & $-1/y$ (empties in $y$ hours) \\
\bottomrule
\end{tabular}
\end{formulabox}

\textbf{Both pipes open together:}
\[
\text{Net rate} = \frac{1}{x} - \frac{1}{y}, \qquad \text{Time} = \frac{1}{\text{Net Rate}}
\]

\begin{examplebox}[Worked Example]
Fill pipe = 8 hr, Empty pipe = 12 hr.
\[
\text{Net rate} = \frac{1}{8} - \frac{1}{12} = \frac{3}{24} - \frac{2}{24} = \frac{1}{24}
\Rightarrow \text{Time} = 24 \text{ hours}
\]
\end{examplebox}

\begin{warnbox}[Common Mistake]
Adding the empty-pipe rate instead of subtracting it. Always treat outlet pipes as \textbf{negative}
contributions to the net rate.
\end{warnbox}

\section{Worker-Day Concept}
Another foundational identity, used constantly for workers joining/leaving problems:
\begin{formulabox}[Worker-Day Identity]
\[
\text{Total Work} = \text{Workers} \times \text{Days}
\]
(only valid when every worker has equal, constant efficiency)
\end{formulabox}

\begin{examplebox}[Worked Example]
20 workers finish a job in 15 days $\Rightarrow$ Total work $= 20 \times 15 = 300$ worker-days.
\end{examplebox}

This ``worker-days'' total is treated as a fixed, conserved quantity --- exactly like total work $=
1$ in the fractional approach, just scaled up to avoid fractions. See
\texttt{02-fast-tricks-and-shortcuts.tex} \S3 for how this is applied to workers-leaving and
workers-joining problems.

\section{Wages Are Proportional to Work --- Not Time}
\begin{warnbox}[Very Common Trap]
\textbf{Money} $\propto$ \textbf{Work Done}, never time spent. If efficiency ratio is $3:5$, wages
split in the ratio $3:5$ --- \textbf{regardless of how many days each person actually worked}, as
long as they worked together on the same job. This trips people up constantly because it feels like
it should depend on time, but wages track \emph{output}, not \emph{hours logged}.
\end{warnbox}

\section{Formula Sheet --- Quick Reference}
\begin{formulabox}[Master Table]
\centering
\begin{tabular}{ll}
\toprule
\textbf{Situation} & \textbf{Formula} \\
\midrule
Work & Rate $\times$ Time \\
Rate (one day's/hour's work) & Work / Time $= 1/\text{Time}$ \\
Time & $1/\text{Rate}$ \\
Combined work (A + B) & $xy/(x+y)$ \\
Combined work ($n$ workers) & Sum of all individual rates, then invert \\
Difference formula & $1/B = 1/(A{+}B \text{ time}) - 1/(A \text{ time})$ \\
Pipe filling & $+1/x$ \\
Pipe emptying & $-1/y$ \\
Efficiency & $1/\text{Time}$ (inversely proportional to time) \\
Wages & $\propto$ Work done, not time \\
Worker-days & Workers $\times$ Days (constant total, equal efficiency) \\
\bottomrule
\end{tabular}
\end{formulabox}

\section{Key Takeaways}
\begin{itemize}[leftmargin=1.4em]
  \item Everything reduces to \textbf{Work = Rate $\times$ Time}, with total work usually normalized
    to 1.
  \item Time and rate are always \textbf{reciprocals}; efficiency and time are always
    \textbf{inversely proportional}.
  \item For ratio-based questions (``A is twice as efficient as B''), assign \textbf{rates directly}
    --- don't guess at days.
  \item Wages split by \textbf{work contributed}, never by time spent.
  \item The worker-day identity (Workers $\times$ Days = constant) is the backbone of every
    joining/leaving problem --- see \texttt{02-fast-tricks-and-shortcuts.tex}.
\end{itemize}

\end{document}
